% Vietnamese translation for OI Public Library
% Source-PDF: 比赛 CONTEST/2021“MINIEYE杯”中国大学生算法设计超级联赛/2021“MINIEYE杯”中国大学生算法设计超级联赛（9）/2021“MINIEYE杯”中国大学生算法设计超级联赛（9）-题解.pdf
\documentclass[11pt]{article}
\usepackage{oipl-vn}

\title{Lời giải HDU Multi-University Training Contest 9 năm 2021}
\author{}
\date{}

\begin{document}
\maketitle

\origtitle{2021 HDU Multi-University Training Contest 9}
\sourcepath{contest/hdu-2021-training-09-solutions.pdf}

\section{Problem A. NJU Emulator}

Trước hết xét ý tưởng nhân đôi. Với mọi $N\ge 2$, ta có thể thu được $N$ từ
$1$ bằng các thao tác $(\times 2)$ hoặc $(\times 2+1)$ trong
$\lfloor\log_2 N\rfloor$ bước. Từ đó có một thuật toán dựng mọi $N$ bằng
$O(\log N)$ lệnh s86 như sau.

Dùng ba lệnh \texttt{p1}, \texttt{dup}, \texttt{add 1} để đưa $1$ và $2$ vào
ngăn xếp, rồi dùng \texttt{p1} cùng các thao tác $(\times 2)$ hoặc
$(\times 2+1)$ tổng cộng $\lfloor\log_2 N\rfloor$ lần để thu được $N$.

Cách này cần nhiều nhất $2\lfloor\log_2 N\rfloor+5$ lệnh để dựng $N$, vẫn còn
quá lớn. Để cải thiện, dùng phương pháp Four Russians: thay cơ số $2$ bằng
một cơ số lớn hơn $k$.

Với mọi $k\ge 2$ và $N\ge 2$, ta có thể thu được $N$ bằng các thao tác
$(\times k+p)$, $0\le p<k$, từ một giá trị $x$ nào đó với $1\le x<k$, trong
$\lfloor\log_k N\rfloor$ bước. Vì thế có thể dùng $2k-1$ lệnh dạng
\texttt{p1}, \texttt{dup}, \texttt{add 1}, \texttt{dup}, \texttt{add 1},
\ldots{} để đưa $1,2,\ldots,k$ vào ngăn xếp, sau đó dùng \texttt{p1},
\texttt{add (x-1)} và $\lfloor\log_k N\rfloor$ thao tác dạng
$(\times k+p)$ để thu được $N$.

Số lệnh khi đó nhiều nhất là $2\lfloor\log_k N\rfloor+2k+2$. Lấy
$k=7$ hoặc $8,9,10,11,12$, ta thu được thuật toán dựng mọi $N<2^{64}$ trong
không quá $60$ lệnh s86, khá gần mục tiêu.

Để cải thiện thêm, chú ý rằng với một cơ số $k$, nếu trong ngăn xếp đã có
$1,2,\ldots,\lfloor k/2\rfloor$ và $k$, thì với mọi $x\le k$ ta có thể lấy
$x$ trong $2$ bước:
\begin{itemize}
  \item dùng \texttt{p1}, \texttt{add x-1} nếu
  $x\le\lfloor k/2\rfloor$;
  \item dùng \texttt{dup}, \texttt{sub k-x} nếu
  $x>\lfloor k/2\rfloor$.
\end{itemize}

Khi đó với $N>k$, ta có thể thu được $N$ từ một $x$ nào đó,
$1\le x\le k$, bằng các thao tác $(\times k+p)$ với
$0\le p\le\lfloor k/2\rfloor$, hoặc $(\times k-p)$ với
$1\le p\le\lfloor k/2\rfloor$, trong $\lfloor\log_k N\rfloor$ bước. Vì vậy
chỉ cần $k+1$ lệnh để đưa $1,2,\ldots,\lfloor k/2\rfloor$ và $k$ vào ngăn
xếp, dựng $x$ trong $2$ bước, rồi áp dụng $\lfloor\log_k N\rfloor$ thao tác
trên để thu được $N$.

Số lệnh nhiều nhất trở thành $2\lfloor\log_k N\rfloor+k+4$. Lấy
$k=16$ hoặc $12,14$, ta có thuật toán dựng mọi $N<2^{64}$ trong không quá
$50$ lệnh s86.

\section{Problem B. Just another board game}

Trước hết xét chiến lược tối ưu nếu không cho phép thao tác ``kết thúc ngay''.
Ta có thể diễn giải lại thao tác của hai người chơi như sau.

Bắt đầu với $r=1$ và $c=1$. Ở mỗi lượt, nếu đến lượt Roundgod, anh ta có thể
đặt $c$ thành một giá trị bất kỳ thỏa $1\le c\le m$. Nếu đến lượt kimoyami,
anh ta có thể đặt $r$ thành một giá trị bất kỳ thỏa $1\le r\le n$. Sau $k$
vòng, trò chơi kết thúc và giá trị trò chơi là $a_{r,c}$.

Xét quá trình theo chiều ngược. Nếu Roundgod đi ở lượt cuối, hiển nhiên anh ta
nên chọn $c$ là vị trí phần tử lớn nhất trên hàng hiện tại; do đó kimoyami,
người đi áp chót, nên chọn hàng có giá trị lớn nhất theo hàng là nhỏ nhất.
Tương tự, nếu kimoyami đi ở lượt cuối, anh ta nên chọn $r$ là vị trí phần tử
nhỏ nhất trên cột hiện tại; do đó Roundgod, người đi áp chót, nên chọn cột có
giá trị nhỏ nhất theo cột là lớn nhất.

Còn một trường hợp đặc biệt $k=1$, tức trò chơi kéo dài nhiều nhất một lượt.
Khi đó rõ ràng Roundgod nên chọn phần tử lớn nhất trên hàng đầu tiên.

Bây giờ xét trường hợp cho phép thao tác ``kết thúc ngay''. Khẳng định là
thao tác này hầu như không thay đổi chiến lược của hai người, ngoại trừ nước
đầu tiên của Roundgod.

Lý do là: nếu sau khi đối thủ vừa đi mà ta chọn thao tác thứ hai, kết quả
không bao giờ tốt hơn việc không làm gì, tức giữ quân cờ đứng yên. Trong
chiến lược tối ưu, đối thủ cũng sẽ không làm gì; hoặc nếu đối thủ cũng chọn
thao tác thứ hai thì kết quả vẫn như vậy. Nếu không, đối thủ chọn một $r$
hoặc $c$ khác với nước vừa đi, điều này rõ ràng không tối ưu, và ta vẫn đạt
được cùng giá trị mà không cần dùng thao tác thứ hai.

Tóm lại, đặt $\mathit{rowmax}_i$ là phần tử lớn nhất trên hàng thứ $i$, và
$\mathit{colmin}_i$ là phần tử nhỏ nhất trên cột thứ $i$. Đáp án được tính
như sau:
\begin{itemize}
  \item nếu $k=1$, đáp án là $\mathit{rowmax}_1$;
  \item nếu $k>1$ và $k$ lẻ, đáp án là
  \[
    \max\left(a_{1,1}, \min_{i=1}^{n}\mathit{rowmax}_i\right);
  \]
  \item nếu $k$ chẵn, đáp án là
  \[
    \max\left(a_{1,1}, \max_{i=1}^{m}\mathit{colmin}_i\right).
  \]
\end{itemize}

Độ phức tạp thời gian tổng thể là $O(nm)$.

\section{Problem C. Dota2 Pro Circuit}

Trước hết xét cách tính $\mathit{worst}_i$ cho mỗi đội; $\mathit{best}_i$
có thể tính tương tự.

Với $i$ và $k$ cố định, để kiểm tra có tồn tại một phương án sao cho có $k$
đội đạt điểm cao hơn đội $i$ hay không, ta cho đội $i$ nhận $b_n$ điểm trong
giải, và chọn $k$ đội có điểm ban đầu cao nhất, trừ đội $i$, nhận lần lượt
$b_1,b_2,\ldots,b_k$ điểm trong giải. Ghép điểm ban đầu và điểm trong giải
theo cách tham lam: đội có điểm ban đầu cao nhất nhận $b_k$ điểm trong giải,
đội có điểm ban đầu cao thứ hai nhận $b_{k-1}$ điểm, và cứ như vậy. Sau đó
kiểm tra xem mọi đội trong $k$ đội được chọn có điểm cuối cùng lớn hơn hẳn
đội $i$ hay không. Phép kiểm tra này tốn $O(n)$ thời gian.

Nếu dùng tìm kiếm nhị phân để tính đáp án cho mọi $i$, độ phức tạp là
$O(n^2\log n)$, quá chậm. Quan sát đơn giản là nếu sắp xếp các đội theo điểm
ban đầu, thì $\mathit{worst}_i$ không tăng. Vì vậy có thể dùng hai con trỏ để
đạt tổng độ phức tạp $O(n^2)$.

\section{Problem D. Into the Woods}

Với một điểm $P(x,y)$, khoảng cách Manhattan từ $P$ đến $O(0,0)$ là
\[
  |x|+|y|=\max\{|x+y|, |x-y|\}.
\]
Giả sử Roundgod đi qua đường đi
$(x_1,y_1),(x_2,y_2),\ldots,(x_n,y_n)$. Gọi $A$ là khoảng cách Manhattan lớn
nhất từ một điểm bất kỳ trên đường đi đến $(0,0)$, khi đó
\[
  A=\max\left\{\max_{i=1}^{n}|x_i+y_i|,
               \max_{i=1}^{n}|x_i-y_i|\right\}.
\]
Theo tính tuyến tính của kỳ vọng, giá trị cần tính là
\[
  E(A)=\sum_{k=1}^{n}\Pr[A\ge k]
      =n-\sum_{k=1}^{n}\Pr[A\le k].
\]

Một mẹo rất hữu ích và đẹp cho loại bước đi ngẫu nhiên hai chiều này là định
nghĩa các biến ngẫu nhiên $a=x+y$ và $b=x-y$. Khi đó $a$ và $b$ độc lập và có
cùng phân phối. Cụ thể, nếu vị trí hiện tại của Roundgod là $(x,y)$, thì ở
mỗi bước, $x+y$ và $x-y$ độc lập tăng hoặc giảm $1$ với xác suất
$\frac12$.

Bài toán được đưa về dạng sau. Ta đi ngẫu nhiên một chiều trên trục $x$ trong
$n$ bước, bắt đầu từ $x=0$; mỗi bước đặt ngẫu nhiên $x=x+1$ hoặc $x=x-1$ với
xác suất $\frac12$. Gọi $p_k$ là xác suất toàn bộ đường đi nằm trong đoạn
$[-k,k]$. Cần tính $p_k$ với $k=1,2,\ldots,n-1$. Khi đó đáp án là
\[
  E(A)=n-\sum_{k=1}^{n-1} p_k^2.
\]

Ta tiếp tục biến đổi sang một bài toán tương đương: xuất phát từ $(0,0)$,
đi $n$ bước, mỗi bước chỉ được đi sang phải hoặc đi lên. Gọi $c_k$ là số
đường đi không chạm đường thẳng $y=x+(k+1)$ và cũng không chạm đường thẳng
$y=x-(k+1)$. Khi đó
\[
  p_k=\frac{c_k}{2^n}.
\]
Để tính số này, cần giải bài toán con sau: cho điểm $P$, đường thẳng
$l_0:y=x+a$ và $l_1:y=x-b$, có bao nhiêu đường đi từ $(0,0)$ đến $P$ mà
không vượt qua $l_0$ hoặc $l_1$? Có thể liên hệ với một bài toán đơn giản
hơn, còn gọi là mở rộng của định lý bỏ phiếu Bertrand: cho điểm $P$ và đường
thẳng $l:y=x+a$, có bao nhiêu đường đi từ $(0,0)$ đến $P$ mà không vượt qua
$l$? Khi $a=0$, số này chính là số Catalan.

Ý tưởng then chốt là nguyên lý phản xạ Andre. Theo nguyên lý phản xạ, khi chỉ
có một đường giới hạn, đáp án là hiệu của hai hệ số nhị thức. Với phiên bản
bao hàm - loại trừ của nguyên lý phản xạ, gọi $f(P)$ là ảnh phản xạ của $P$
qua $l_0$, $g(P)$ là ảnh phản xạ của $P$ qua $l_1$, và $C(P)$ là số đường đi
từ $(0,0)$ đến $P$, khi đó đáp án là
\[
  C(P)-C(f(P))-C(g(P))+C(g(f(P)))+C(f(g(P)))-\cdots .
\]

Suy ra có thể tính $c_k$ bằng
\[
  c_k=
  \sum_{i=L}^{R}\binom{n}{i}
  +2\sum_{j\ge 1}(-1)^j\sum_{i=L}^{R}
    \binom{n}{i-(k+1)j},
\]
trong đó
\[
  L=\left\lceil\frac{n-k}{2}\right\rceil,\qquad
  R=\left\lfloor\frac{n-k}{2}\right\rfloor.
\]

Nếu tiền xử lý tổng tiền tố của $\binom{n}{i}$, thì $c_k$ tính được trong
$O(n/k)$ thời gian. Do đó độ phức tạp tổng thể là
\[
  \sum_{i=1}^{n}\frac{n}{i}=O(n\log n).
\]

\section{Problem E. Did I miss the lethal?}

Bài này có thể diễn giải lại thành một trò chơi.

Ta có $n$ lá bài trên tay; lá thứ $i$ có hai giá trị tương ứng là $d_i$ và
$a_i$. Ta và đối thủ, tức người chia bài trong Hearthstone, luân phiên đi ở
mỗi vòng:
\begin{enumerate}
  \item Ta chọn một lá có $d_i$ và $a_i$, bỏ lá đó và gây $d_i$ sát thương.
  \item Đối thủ chọn $a_i$ lá và bỏ chúng.
\end{enumerate}
Ta muốn tối đa hóa tổng sát thương gây ra, còn đối thủ muốn tối thiểu hóa nó.
Cần tính đáp án khi cả hai chơi tối ưu.

Chú ý rằng mỗi $a_i$ thuộc $\{1,2,3,4\}$, nên ta chia toàn bộ $n$ lá thành
$4$ đống: mọi lá có $a_i=j$ được đặt vào đống thứ $j$, với $1\le j\le 4$.
Gọi $n_1,n_2,n_3,n_4$ là số lá ban đầu trong từng đống. Rõ ràng có một chiến
lược tham lam: một khi ta hoặc đối thủ quyết định bỏ một lá khỏi một đống,
lá có sát thương lớn nhất trong đống đó sẽ được chọn. Vì vậy với mỗi đống,
sắp xếp các lá theo $d_i$, rồi giải bằng quy hoạch động.

Đặt $dp[x_1][x_2][x_3][x_4][0]$ là đáp án khi đống thứ $j$ còn $x_j$ lá với
mọi $1\le j\le 4$, và đến lượt ta chọn lá. Đặt
$dp[x_1][x_2][x_3][x_4][k]$, với $k=1,2,3,4$, là đáp án khi đống thứ $j$ còn
$x_j$ lá và đến lượt đối thủ chọn $k$ lá. Các chuyển trạng thái đều tính được
trong $O(1)$.

Tổng số trạng thái là $5n_1n_2n_3n_4$, với
\[
  \sum_{i=1}^{4} n_i\le 200.
\]
Số trạng thái đạt lớn nhất khi
$n_1=n_2=n_3=n_4=50$, tức $50^4\cdot 5$, đủ nhỏ để dễ dàng qua giới hạn thời
gian.

\section{Problem F. Guess The Weight}

Trước hết xét chiến lược tối ưu của uuzlovetree, phần này khá đơn giản: sau
khi rút lá đầu tiên, anh ta nhìn các lá còn lại trong bộ bài rồi chọn
``Small'' hoặc ``Large'' tùy phía nào, so với lá đầu tiên đã rút, có nhiều lá
hơn.

Quan sát then chốt là với mọi bộ bài, tồn tại một ``điểm giữa'' $\mathit{mp}$
sao cho trong chiến lược tối ưu, uuzlovetree nên chọn ``Large'' nếu lá đầu
tiên có số không vượt quá $\mathit{mp}$, và nên chọn ``Small'' nếu lá đầu
tiên có số lớn hơn $\mathit{mp}$.

Giả sử đa tập lá bài hiện tại là $SC$, và số lá trong bộ là $s$, tức
$|SC|=s$. Đặt
\[
  L=\{x\mid x\in SC,\ x\le \mathit{mp}\},\qquad
  R=\{x\mid x\in SC,\ x>\mathit{mp}\}.
\]
Khi xét mọi khả năng của hai lá đầu tiên trên bộ bài, xác suất rút được lá
thứ hai là
\[
  \Pr[\mathrm{DrawSecond}]
  =1-\frac{
    \sum_{x\in L,y\in L}[x\ge y]
    +\sum_{x\in R,y\in R}[x\le y]
  }{s(s+1)}.
\]

Do đó, nếu dùng cây đoạn để duy trì số lá và số cặp lá có giá trị khác nhau
trong mỗi đoạn, ta có thể tính xác suất này trong $O(\log\max a_i)$. Điểm
giữa cũng tìm được trong cùng độ phức tạp bằng cách đi xuống trên cây đoạn.
Độ phức tạp thời gian tổng thể là
\[
  O(n+\max a_i+q\log\max a_i).
\]

\section{Problem G. Boring data structure problem}

Duy trì phía trái, tức phần bên trái phần tử ở giữa, và phía phải, tức phần
bên phải phần tử ở giữa, bằng hai danh sách liên kết đôi $L$ và $R$. Khi đó
hai thao tác chèn đầu tiên đều thực hiện được trong $O(1)$ thời gian cho mỗi
truy vấn.

Sau khi thực hiện một thao tác, điều chỉnh kích thước hai danh sách liên kết
đôi như sau:
\begin{itemize}
  \item nếu $|L|>|R|$, chuyển phần tử cuối của $L$ sang đầu $R$;
  \item nếu $|L|<|R|+1$, chuyển phần tử đầu của $R$ sang cuối $L$.
\end{itemize}
Việc điều chỉnh này hiển nhiên có độ phức tạp khấu hao $O(1)$ cho mỗi truy
vấn.

Khi xóa phần tử, chỉ cần đánh dấu phần tử đó là đã xóa và giảm kích thước
danh sách liên kết đôi tương ứng đi một. Đây là thao tác ``lười''; ta chỉ lan
truyền việc xóa khi cần thực hiện thao tác trên phần tử đó, chẳng hạn di
chuyển hoặc truy vấn. Vì thế thao tác xóa cũng có độ phức tạp khấu hao
$O(1)$.

Do đó độ phức tạp thời gian tổng thể là $O(n+q)$.

\section{Problem H. Integers Have Friends 2.0}

Trước hết quan sát rằng nếu lấy $m=2$ và chọn toàn bộ $a_i$ lẻ hoặc toàn bộ
$a_i$ chẵn, ta luôn bảo đảm có đáp án ít nhất $\lceil n/2\rceil$.

Chọn ngẫu nhiên hai chỉ số phân biệt $1\le x<y\le n$. Xác suất để đáp án tối
ưu chứa cả hai chỉ số này ít nhất là $\frac14$. Nếu cả hai cùng nằm trong đáp
án, thì $m$ phải là một ước của $|a_x-a_y|$. Thực ra ta chỉ cần kiểm tra các
thừa số nguyên tố của $|a_x-a_y|$. Mọi số dương $x$ thỏa
$x\le 4\cdot 10^{12}$ có nhiều nhất $11$ thừa số nguyên tố phân biệt, và việc
tìm tất cả chúng có thể làm dễ dàng trong $O(\sqrt{\max a_i})$ nếu tiền xử
lý mọi số nguyên tố không vượt quá $2\cdot 10^6$.

Lặp quá trình trên $K$ lần. Xác suất sai nhiều nhất là $(3/4)^K$. Thời gian
chạy là
\[
  O(K\sqrt{\max a_i}+11Kn).
\]
Lấy $K=40$ cho xác suất lỗi không đáng kể và vẫn dễ dàng nằm trong giới hạn
thời gian.

\section{Problem I. Little Prince and the garden of roses}

Rõ ràng có thể giải bài toán độc lập cho từng màu hoa hồng. Cố định một màu
$c$. Ta cần gán màu cho cuống của các hoa hồng có màu $c$. Một phép quy về
quen thuộc trên lưới như sau: xem lưới $n\times n$ là một đồ thị hai phía
$G=(V=L\cup R,E)$ với $|L|=|R|=n$, và xem mỗi hoa hồng ở vị trí $(u,v)$ là
một cạnh nối đỉnh thứ $u$ ở phía trái với đỉnh thứ $v$ ở phía phải. Khi đó
bài toán trở thành tô màu cạnh của đồ thị hai phía: gán màu cho mỗi cạnh sao
cho không có hai cạnh kề cùng màu, đồng thời tối thiểu hóa số màu khác nhau
được dùng.

Vậy cần bao nhiêu màu? Gọi $\Delta$ là bậc lớn nhất của một đỉnh trong đồ thị
hai phía. Hiển nhiên cần ít nhất $\Delta$ màu. Ta có thể chứng minh rằng
$\Delta$ màu cũng đủ với đồ thị hai phía; đây là trường hợp đặc biệt của định
lý Vizing, phát biểu rằng chỉ số tô màu cạnh của đồ thị tổng quát là
$\Delta$ hoặc $\Delta+1$. Có hai cách làm.

Cách thứ nhất là một thuật toán xây dựng chạy trong $O(|E||V|)$. Tô màu các
cạnh lần lượt theo một thứ tự nào đó. Gọi $(x,y)$ là cạnh hiện tại. Nếu tồn
tại màu $c$ đang rảnh ở cả đỉnh $x$ lẫn đỉnh $y$, ta tô trực tiếp cạnh
$(x,y)$ bằng $c$. Nếu không có màu như vậy, tồn tại một cặp màu $c_1,c_2$ sao
cho $c_1$ có ở $x$ nhưng không có ở $y$, còn $c_2$ có ở $y$ nhưng không có ở
$x$. Ta làm cho đỉnh $y$ rảnh màu $c_2$. Gọi $z$ là đầu còn lại của cạnh kề
$y$ có màu $c_2$. Nếu $z$ rảnh màu $c_1$, ta tô cạnh $(x,y)$ bằng $c_2$ và
đổi màu cạnh $(y,z)$ sang $c_1$. Như vậy ta thực hiện một phép hoán đổi theo
dây xen kẽ. Nếu $z$ không rảnh màu $c_1$, gọi $w$ là đầu còn lại của cạnh kề
$z$ có màu $c_1$. Nếu $w$ rảnh màu $c_2$, lại có thể hoán đổi. Vì đồ thị là
hai phía, cuối cùng ta sẽ tìm được một dây xen kẽ. Dùng DFS để tìm dây này;
mỗi dây chứa không quá $n$ đỉnh, nên độ phức tạp tổng thể là $O(|E||V|)$.

Cách thứ hai dễ hiểu hơn nhưng chậm hơn. Trước hết thêm các cạnh giả giữa các
đỉnh trong $L$ và $R$, có thể có bội cạnh, sao cho mọi đỉnh ở $L$ và $R$ đều
có bậc $\Delta$. Sau đó tìm một ghép hoàn hảo bất kỳ, tô toàn bộ cạnh trong
ghép bằng cùng một màu, xóa các cạnh của ghép, rồi lặp lại $\Delta$ lần. Cách
này cũng cho một phép tô cạnh bằng $\Delta$ màu khác nhau. Tính đúng đắn có
thể chứng minh bằng định lý Hall, ở đây lược bỏ. Độ phức tạp thời gian là
\[
  O(\Delta |E|\sqrt{|V|}),
\]
và vẫn có thể được chấp nhận nếu cài đặt tốt.

Quy về theo $n$, độ phức tạp tổng thể là $O(n^3)$ cho cách thứ nhất và
$O(n^{3.5})$ cho cách thứ hai.

\section{Problem J. Unfair contest}

Sau khi sắp xếp $a_1,a_2,\ldots,a_{n-1}$ và
$b_1,b_2,\ldots,b_{n-1}$, điểm mới thêm $a_n$ hoặc $b_n$ có thể rơi vào ba
loại: nằm trong $s$ điểm cao nhất, nằm trong $t$ điểm thấp nhất, hoặc không
thuộc hai nhóm đó. Với mỗi loại, ta có một khoảng giá trị hợp lệ có thể có
của $a_n$ hoặc $b_n$; phần còn lại chỉ là phân tích trường hợp cẩn thận và
đầy đủ. Cần chú ý các trường hợp $s=0$ hoặc $t=0$.

Độ phức tạp thời gian là $O(n\log n)$ do cần sắp xếp.

\section{Problem K. ZYB's kingdom}

Trước hết xét cách giải khi $k=1$ với mỗi truy vấn. Sự kiện loại $1$ tương
đương với việc: với mỗi cây con được tách ra bởi $v$, gọi $s$ là tổng $c$
trong cây con này; khi đó với mỗi đỉnh $u$ trong cây con, ta cộng $s$ rồi
trừ $c_u$ khỏi thu nhập của nó.

Tồn tại một thuật toán quen thuộc $O(q\sqrt n\log n)$ bằng cách tách các đỉnh
có bậc không quá $\sqrt n$ và ít nhất $\sqrt n$, rồi dùng cây đoạn để cập
nhật đáp án. Tuy nhiên có một cách khéo hơn nhiều, giải được trong
$O(q\log n)$.

Ý tưởng là trước hết chọn gốc tùy ý cho cây rồi thực hiện phân rã nặng nhẹ.
Với mỗi sự kiện loại $1$, ta chỉ cập nhật cây con nối với cha của $v$ và cây
con nối với đứa con nặng duy nhất của $v$. Việc này tốn $O(\log n)$ thời
gian, vì chỉ cập nhật hai cây con. Với mỗi sự kiện loại $2$, sau khi truy vấn
đáp án của $v$ trên cây đoạn, phần nào còn thiếu? Ta cần cộng đóng góp của
những cập nhật tại $u$ sao cho $u$ là tổ tiên của $v$, nhưng $v$ không nằm
trong cây con nối với con nặng của $u$. Nhờ tính chất của phân rã nặng nhẹ,
chỉ có nhiều nhất $O(\log n)$ đỉnh như vậy. Vì thế, nếu lưu các cập nhật vào
một mảng khi gặp sự kiện loại $1$, và khi gặp sự kiện loại $2$ thì lần lượt
đi lên theo các chuỗi nặng, ta giải được mỗi sự kiện trong $O(\log n)$. Độ
phức tạp tổng thể là $O(q\log n)$.

Khi $k>1$ thì sao? Về cơ bản mọi thứ vẫn giống vậy; chỉ cần áp dụng cùng ý
tưởng cho từng thành phần. Ở đây có thể cần xây một cấu trúc cây, còn gọi là
cây ảo, theo $k$ đỉnh bị cấm, rồi giải đệ quy trên cấu trúc cây này. Độ phức
tạp tổng thể là
\[
  O\left(\sum k\log n\right).
\]

\end{document}
